\section{Практическая часть}
В листинге~\ref{lst:source} представлен исходный код реализованной программы.
\begin{listing}[Исходный код программы]{lst:source}{Matlab}
clear; clc; close all;

X = [7.76,6.34,5.11,7.62,8.84,4.68,8.65,6.90,8.79,6.61, ...
    6.62,7.13,6.75,7.28,7.74,7.08,5.57,8.20,7.78,7.92, ...
    6.00,4.88,6.75,6.56,7.48,8.51,9.06,6.94,6.93,7.79, ...
    5.71,5.93,6.81,5.76,5.88,7.05,7.22,6.67,5.59,6.57, ...
    7.28,6.22,6.31,5.51,6.69,7.12,7.40,6.86,7.28,6.82, ...
    7.08,7.52,6.81,7.55,4.89,5.48,7.74,5.10,8.17,7.67, ...
    7.07,5.80,6.10,7.15,7.88,9.06,6.85,4.88,6.74,8.76, ...
    8.53,6.72,7.21,7.42,8.29,8.56,9.25,6.63,7.49,6.67, ...
    6.79,5.19,8.20,7.97,8.64,7.36,6.72,5.90,5.53,6.44, ...
    7.35,5.18,8.25,5.68,6.29,6.69,6.08,7.42,7.10,7.14, ...
    7.10,6.60,6.35,5.99,6.17,9.05,6.01,7.77,6.27,5.81, ...
    7.80,9.89,4.39,6.83,6.53,8.15,6.68,6.87,6.31,6.83];

N = length(X);
gamma = 0.9;

mu_final  = mean(X);
var_final = var(X);

n_array = 1:N;
mu_hat  = NaN(1, N);
mu_low  = NaN(1, N);
mu_up   = NaN(1, N);

s2_hat  = NaN(1, N);
var_low = NaN(1, N);
var_up  = NaN(1, N);

for n = 2:N
    xn = X(1:n);
    
    mu_hat(n) = mean(xn);
    s2_hat(n) = var(xn);
    
    alpha_t = (1 + gamma) / 2;
    t_val   = tinv(alpha_t, n - 1);
    
    delta_mu = t_val * sqrt(s2_hat(n) / n);
    mu_low(n) = mu_hat(n) - delta_mu;
    mu_up(n)  = mu_hat(n) + delta_mu;
    
    alpha_chi_L = (1 - gamma) / 2;
    alpha_chi_R = (1 + gamma) / 2;
    
    chi_L = chi2inv(alpha_chi_L, n - 1);
    chi_R = chi2inv(alpha_chi_R, n - 1);
    
    var_low(n) = ((n - 1) * s2_hat(n)) / chi_R;
    var_up(n)  = ((n - 1) * s2_hat(n)) / chi_L;
end

fprintf('--- РЕЗУЛЬТАТЫ ДЛЯ N = %d, gamma = %g ---\n', N, gamma);
fprintf('Математическое ожидание:\n');
fprintf('  Точечная оценка: %.4f\n', mu_hat(N));
fprintf('  Доверительный интервал: [%.4f ; %.4f]\n\n', mu_low(N), mu_up(N));

fprintf('Дисперсия:\n');
fprintf('  Точечная оценка: %.4f\n', s2_hat(N));
fprintf('  Доверительный интервал: [%.4f ; %.4f]\n', var_low(N), var_up(N));


figure('Name', 'Оценки математического ожидания', 'Color', 'w');
hold on; grid on;

plot([1, N], [mu_final, mu_final], '--', 'Color', color_final, 'LineWidth', 2, 'DisplayName', '$\hat{\mu}(\vec{x}_N)$');

plot(n_array, mu_hat, '-o', 'LineWidth', 1.5, 'MarkerSize', 4, 'DisplayName', '$\hat{\mu}(\vec{x}_n)$');
plot(n_array, mu_up, '-', 'LineWidth', 1.5, 'DisplayName', '$\overline{\mu}(\vec{x}_n)$');
plot(n_array, mu_low, '-', 'LineWidth', 1.5, 'DisplayName', '$\underline{\mu}(\vec{x}_n)$');

title(['Доверительный интервал для матожидания (\gamma = ', num2str(gamma), ')']);
xlabel('Объем выборки n');
ylabel('Значение y');
legend('Location', 'best', 'Interpreter', 'latex', 'FontSize', 14);
xlim([1 N]);

figure('Name', 'Оценки дисперсии', 'Color', 'w');
hold on; grid on;

plot([1, N], [var_final, var_final], '--', 'Color', color_final, 'LineWidth', 2, 'DisplayName', '$S^2(\vec{x}_N)$');

plot(n_array, s2_hat, '-o', 'LineWidth', 1.5, 'MarkerSize', 4, 'DisplayName', '$S^2(\vec{x}_n)$');
plot(n_array, var_up, '-', 'LineWidth', 1.5, 'DisplayName', '$\overline{\sigma^2}(\vec{x}_n)$');
plot(n_array, var_low, '-', 'LineWidth', 1.5, 'DisplayName', '$\underline{\sigma^2}(\vec{x}_n)$');

title(['Доверительный интервал для дисперсии (\gamma = ', num2str(gamma), ')']);
xlabel('Объем выборки n');
ylabel('Значение z');
legend('Location', 'best', 'Interpreter', 'latex', 'FontSize', 14);
xlim([1 N]);
\end{listing}

В листинге~\ref{lst:calc} приведены результаты расчетов указанных величин.
\begin{listing}[Результаты расчетов указанных величин]{lst:calc}{Matlab}
--- РЕЗУЛЬТАТЫ ДЛЯ N = 120, gamma = 0.9 ---
Математическое ожидание:
  Точечная оценка: 6.9445
  Доверительный интервал: [6.7807 ; 7.1083]

Дисперсия:
  Точечная оценка: 1.1720
  Доверительный интервал: [0.9588 ; 1.4710]
\end{listing}

На рисунках~\ref{img:m-graph-1.pdf}-\ref{img:m-graph-2.pdf} представлены 
график прямой \( y=\hat{\mu}(\vec{x}_N) \),
графики функций \( y=\hat{\mu}(\vec{x}_n) \), \( y=\underline{\mu}(\vec{x}_n) \),
\( y=\overline{\mu}(\vec{x}_n) \), как функций объема выборки \( n \), 
где \( n \) принимает значения от 1 до \( N \).

\jimg{m-graph-1.pdf}{График \( \gamma \)-доверительного интервал для математического ожидания}
\jimg{m-graph-2.pdf}{Увеличенный график \( \gamma \)-доверительного интервал для математического ожидания}

На рисунках~\ref{img:d-graph-1.pdf}-\ref{img:d-graph-2.pdf} представлены 
график прямой \( z=S^2(\vec{x}_N) \),
графики функций \( z=S^2(\vec{x}_n) \), \( z=\underline{\sigma^2}(\vec{x}_n) \),
\( z=\overline{\sigma^2}(\vec{x}_n) \), как функций объема выборки \( n \),
где \( n \) принимает значения от 1 до \( N \).

\jimg{d-graph-1.pdf}{График \( \gamma \)-доверительного интервал для дисперсии}
\jimg{d-graph-2.pdf}{Увеличенный график \( \gamma \)-доверительного интервал для дисперсии}
